\documentclass[10pt,a4paper]{amsart}
\usepackage[margin=19mm]{geometry}
\usepackage{amsmath,amssymb,mathtools}
\usepackage[scr=rsfs,cal=euler]{mathalfa}
\usepackage{xfrac}
\usepackage[unicode,hidelinks,bookmarks=false]{hyperref}
\newcommand{\ie}{i.e.\ }
\newcommand{\cf}{cf.\ }
\newcommand{\resp}{respectively\ }
\newcommand{\Subsection}{\subsection}
\providecommand{\nobreakdash}{\nobreak\hbox{-}}
\setlength{\emergencystretch}{2em}
\multlinegap0pt
\newtheorem{theorem}{Theorem}[section]
\newtheorem{definition}[theorem]{Definition}
\newtheorem{lemma}[theorem]{Lemma}
\newtheorem{corollary}[theorem]{Corollary}
\newtheorem{remark}[theorem]{Remark}
\newtheorem{example}[theorem]{Example}
\newcommand{\C}{\mathbb{C}}
\newcommand{\R}{\mathbb{R}}
\newcommand{\Z}{\mathbb{Z}}
\newcommand{\T}{\mathbb{T}}
\newcommand{\Hilb}{\mathcal{H}}
\newcommand{\Pic}{\operatorname{Pic}}
\newcommand{\Ker}{\operatorname{Ker}}
\begin{document}
\title[Quantization of Contact 3-Manifolds and the Reeb Gravitation]{Quantization of Contact 3-Manifolds and the Reeb Gravitational Field}
\author{Ali M. Elgindi}
\date{}
\maketitle
\noindent\emph{Source and licence.} Quantization of Contact 3-Manifolds and the Reeb Gravitational Field. Version arXiv:2606.16529v1 (CC BY 4.0). This material is distributed under the Creative Commons Attribution 4.0 International licence, http://creativecommons.org/licenses/by/4.0/, as recorded in the arXiv OAI licence metadata for this exact version and shown on the abstract page https://arxiv.org/abs/2606.16529v1. Full rights notice: licenses/arxiv-2606.16529-CC-BY-4.0.txt.\par\smallskip
\noindent\textit{Editorial scope.} Counted unit: the original \texttt{theorem} environment \texttt{-} at source lines 88--94 together with its complete proof at lines 96--102; the delivered document keeps source lines 76--103 so that every referenced label and every citation resolves inside it. Native editable source is the paper's own concatenated TeX; the AMS wrapper is editorial formatting only. No publisher class, upstream script or unrelated artwork is redistributed.
\section{The Quantum Hilbert Space}

On the quantum locus $L$, the contact structure $\xi|_L$ is a complex line bundle. Let $\kappa_{\xi|_L} = \wedge^1 (\xi|_L)^*$ be the canonical bundle of the polarization. Since $L$ is a disjoint union of circles, $H^2(L;\Z) = 0$, so a unique square root $\kappa_{\xi|_L}^{1/2}$ exists.  (see [2], [12])

\begin{definition}[Quantum Hilbert Space]
The quantum Hilbert space associated to $(M,\xi)$ is
\[
\Hilb_\xi = H^0(L, L_\xi|_L \otimes \kappa_{\xi|_L}^{1/2}),
\]
where $H^0$ denotes holomorphic sections.
\end{definition}


\begin{theorem}
$\Hilb_\xi$ is finite-dimensional. If $L = \bigsqcup_i L_i$ are the connected components of complex tangents, then:
\[
\dim \Hilb_\xi = \sum_i \deg(L_\xi|_{L_i}),
\]
where $\deg$ is the first Chern number.
\end{theorem}

\begin{proof}
Each $L_i$ is a compact complex curve of genus $0$ (a Riemann sphere with boundary). By the Riemann-Roch theorem,
\[
\dim H^0(L_i, L_\xi|_{L_i} \otimes \kappa^{1/2}) = \deg(L_\xi|_{L_i}) + \frac{1}{2}\deg(\kappa) + 1 - g_i.
\]
Since $\deg(\kappa) = -2$ and $g_i = 0$, we obtain $\dim = \deg(L_\xi|_{L_i})$. Summing over $i$ gives the result.
\end{proof}


\end{document}
